\documentclass[12pt]{article}
\usepackage[utf8]{inputenc}
\usepackage[T1]{fontenc}
\usepackage{hyperref}
\title{Telefol Language Resources for K-12 Curriculum Development}
\author{Compiled by Claude Sonnet 5 (claude-sonnet-5) for Andr\'as Kornai}
\date{}
\begin{document}
\maketitle

\section{Introduction}
Telefol (ISO 639-3 tlf) is a Trans-New Guinea language of Papua New Guinea. Estimates of its speakers range from about 4,500 to 5,400. Sources describe two main dialects, Southern (Feramin) and Northern (Telefolmin). The search found little material for schools. Nothing found is a primer, reader or curriculum.

\section{Primary Educational Resources}
No Telefol primers, readers, textbooks or literacy programme materials turned up in the searches. This gap is real for this search and does not show that none exist.

\section{Linguistic and Reference Materials}
\begin{itemize}
\item Healey, P. and A. \emph{Telefol Dictionary} (Pacific Linguistics, 1977), open access at the Australian National University: \url{https://openresearch-repository.anu.edu.au/items/e2c324d2-d8c7-455c-98de-1cc204051024}. It is based mainly on the Kialikmin parish dialect of the Northern dialect.
\item Healey's 1964--1965 studies of Telefol phonology, noun phrases, clause structure and verb phrases are listed in the WALS reference database: \url{https://wals.info/languoid/lect/wals_code_tlf}.
\item Overview article: \url{https://en.wikipedia.org/wiki/Telefol_language}.
\end{itemize}

\section{Religious and Community Resources}
Global Recordings Network lists Telefol and notes that a New Testament and portions of scripture exist: \url{https://globalrecordings.net/en/language/tlf}. The search results do not explain the page's literacy figure. I did not verify it.

\section{Digital and Audio Resources}
Global Recordings Network (above) lists the language. The catalogue entry at the Endangered Languages Project is \url{https://endangeredlanguages.com/lang/10675}. I found no other digital resources.

\section{Curriculum Development Context}
The only substantial resource found is the 1977 dictionary. A mother-tongue curriculum would have to be built from this dictionary and from direct work with the community. Check SIL and Papua New Guinea literacy sources for further materials.
\end{document}